\section{Hash Maps}

A \textbf{hash map} stores \textbf{(key, value)} pairs and supports fast lookup by key.  
In a hash map, insert/search/remove are \textbf{$O(1)$}. Hash maps do \textbf{not maintain any order}, so the physical placement of keys in the table has no relation to insertion order or sorted order.

\subsection{Core Idea}
A hash map is implemented using an \textbf{array}. 
A \textbf{hash function} converts a key into an integer index:
\[
\text{index} = h(\text{key}) \bmod m
\]
where $m$ is the table size.

The same key always hashes to the same index.

\subsection{Example Hash Function}
An example of a hash function for strings:
\begin{itemize}
  \item sum ASCII values of characters
  \item mod by table size to stay in bounds
\end{itemize}

\begin{verbatim}
int hash(String key, int tableSize) {
    int sum = 0;
    for (char c : key.toCharArray()) {
        sum += c;
    }
    return sum % tableSize;
}
\end{verbatim}

\subsection{Step-by-Step Insertion Example (Open Addressing)}

Initial table size: $m = 2$ (indices 0 and 1)

\begin{center}
\begin{tabular}{c c}
\toprule
Index & Value \\
\midrule
0 & -- \\
1 & -- \\
\bottomrule
\end{tabular}
\end{center}

\textbf{Insert:} \texttt{("Alice", "NYC")}  
Assume $\text{hash(Alice)} = 25 \Rightarrow 25 \bmod 2 = 1$

\begin{center}
\begin{tabular}{c c}
\toprule
Index & Value \\
\midrule
0 & -- \\
1 & (Alice, NYC) \\
\bottomrule
\end{tabular}
\end{center}

The table is now half full $\Rightarrow$ \textbf{resize and rehash}. We do this to reduce collisions, keeping insert, search, and remove $O(1)$ time.

\subsection{Rehashing (after resize to size 4)}
Recompute all indices using new size $m = 4$.

\[
25 \bmod 4 = 1
\]

\begin{center}
\begin{tabular}{c c}
\toprule
Index & Value \\
\midrule
0 & -- \\
1 & (Alice, NYC) \\
2 & -- \\
3 & -- \\
\bottomrule
\end{tabular}
\end{center}

\textbf{Insert:} \texttt{("Brad", "Chicago")}  
Assume hash = 27, $27 \bmod 4 = 3$

\begin{center}
\begin{tabular}{c c}
\toprule
Index & Value \\
\midrule
0 & -- \\
1 & (Alice, NYC) \\
2 & -- \\
3 & (Brad, Chicago) \\
\bottomrule
\end{tabular}
\end{center}

The table is now half full again $\Rightarrow$ \textbf{resize and rehash}.  

\subsection{Rehashing (after resize to size 8)}
Recompute all indices using new size $m = 8$.

\[
25 \bmod 8 = 1,\quad 27 \bmod 8 = 3
\]

\begin{center}
\begin{tabular}{c c}
\toprule
Index & Value \\
\midrule
0 & -- \\
1 & (Alice, NYC) \\
2 & -- \\
3 & (Brad, Chicago) \\
4 & -- \\
5 & -- \\
6 & -- \\
7 & -- \\
\bottomrule
\end{tabular}
\end{center}

\subsection{Collision Example (Open Addressing)}
\textbf{Insert:} \texttt{("Collin", "Seattle")}  
Assume hash = 33, $33 \bmod 8 = 1$ $\Rightarrow$ collision with Alice. In this example, we handle collisions with \textbf{open addressing}, where we find the next empty slot in the array so that multiple keys can still be stored even when they hash to the same index. Other collision-handling strategies exist, such as \textbf{chaining}, where each index stores a linked list of key–value pairs instead of using probing.


\begin{itemize}
  \item index 1 is occupied
  \item try index 2 (empty)
\end{itemize}

\begin{center}
\begin{tabular}{c c}
\toprule
Index & Value \\
\midrule
0 & -- \\
1 & (Alice, NYC) \\
2 & (Collin, Seattle) \\
3 & (Brad, Chicago) \\
4 & -- \\
5 & -- \\
6 & -- \\
7 & -- \\
\bottomrule
\end{tabular}
\end{center}

\subsection{Search Example (After Collision)}

Because collisions push elements forward, search must also \textbf{probe forward} until the key is found or an empty slot is reached.

\begin{verbatim}
get(key):
    index = h(key) % m
    while table[index] is not empty:
        if table[index].key == key:
            return table[index].value
        index = (index + 1) % m   // linear probing
    return NOT_FOUND
\end{verbatim}

\subsubsection*{Example 1: Key Found}
Search for \texttt{"Collin"}:

\begin{itemize}
  \item hash("Collin") = 33 $\Rightarrow 33 \bmod 8 = 1$
  \item index 1 contains Alice $\Rightarrow$ not a match
  \item probe next index (2)
  \item index 2 contains Collin $\Rightarrow$ found
\end{itemize}

\[
\text{get("Collin")} = \text{"Seattle"}
\]

\subsubsection*{Example 2: Key Not Found}
Search for \texttt{"Daniel"} (assume it also hashes to index 1 for this example):

\begin{itemize}
  \item start at index 1
  \item index 1: Alice (not match)
  \item index 2: Collin (not match)
  \item index 3: Brad (not match)
  \item index 4: empty slot reached $\Rightarrow$ stop
\end{itemize}

\[
\text{get("Daniel")} = \text{NOT\_FOUND}
\]

Reaching an empty slot means the key does not exist, since insertion would have stopped there.

\subsection{Duplicate Key (Overwrite)}
Keys are typically \textbf{unique} (no duplicate keys). Inserting an existing key usually \textbf{overwrites} its value.

Insert \texttt{("Alice", "Boston")}:

\begin{itemize}
  \item hash gives index 1
  \item key already exists
  \item value is overwritten
\end{itemize}

\begin{center}
\begin{tabular}{c c}
\toprule
Index & Value \\
\midrule
0 & -- \\
1 & (Alice, Boston) \\
2 & (Collin, Seattle) \\
3 & (Brad, Chicago) \\
4 & -- \\
5 & -- \\
6 & -- \\
7 & -- \\
\bottomrule
\end{tabular}
\end{center}

\subsection{Deletion + Tombstones (Open Addressing)}
With open addressing, we \textbf{cannot} delete by setting a slot to empty (\texttt{--}), because that can break the probe chain and cause searches to stop early.  
Instead, we mark deleted slots with a \textbf{tombstone} (e.g. \texttt{DEL}).

\subsubsection*{Tombstone Example}
\textbf{Remove:} \texttt{"Brad"} (at index 3) $\Rightarrow$ mark as \texttt{DEL} (tombstone), not empty.

\begin{center}
\begin{tabular}{c c}
\toprule
Index & Value \\
\midrule
0 & -- \\
1 & (Alice, Boston) \\
2 & (Collin, Seattle) \\
3 & DEL \\
4 & -- \\
5 & -- \\
6 & -- \\
7 & -- \\
\bottomrule
\end{tabular}
\end{center}

\textbf{How search changes:}
\begin{itemize}
  \item \texttt{DEL} means ``keep probing'' (key might be later in the chain)
  \item \texttt{--} means ``stop'' (key does not exist)
\end{itemize}

\textbf{How insert changes:}
\begin{itemize}
  \item \texttt{DEL} can be reused as an insertion spot
\end{itemize}

\begin{verbatim}
get(key):
    idx = h(key) % m
    while table[idx] != EMPTY:
        if table[idx] != DEL and table[idx].key == key:
            return table[idx].value
        idx = (idx + 1) % m
    return NOT_FOUND
\end{verbatim}

\subsection{Time Complexity}

\begin{center}
\begin{tabular}{l c c}
\toprule
\textbf{Operation} & \textbf{Average Case} & \textbf{Worst Case} \\
\midrule
Insert & $O(1)$ & $O(n)$ \\
Search & $O(1)$ & $O(n)$ \\
Remove & $O(1)$ & $O(n)$ \\
\bottomrule
\end{tabular}
\end{center}

Worst case happens when many keys collide, forcing long probing/chains, but this is very unlikely with a good hash function and regular resizing.